% GENERATED FILE. Edit canonical lessons, then run npm run generate:books.
% canonical-source-hash: fnv1a64:2912a698279a5122
% canonical-lessons: ML-C521-samtrpti, ML-C521-visadam, ML-C521-puncirikkuka, ML-C521-grhaturatvam, ML-C521-netuvirppituka

\chapter{Smiles and Sighs}
\label{ch:ml-smiles-and-sighs}
\hlchaptermodality{521}

\begin{chapteropening}
\textbf{By the end of this chapter:} \emph{I can say satisfaction, melancholy, to smile, nostalgia and to sigh.}

Complete the last lesson of chapter 521: I can say satisfaction, melancholy, to smile, nostalgia and to sigh.
\end{chapteropening}

\section[\texorpdfstring{\ml{സംതൃപ്തി} (saṁtṛpti)}{സംതൃപ്തി (saṁtṛpti)}]{\ml{സംതൃപ്തി} (saṁtṛpti) — satisfaction; contentment}

\label{lesson:ML-C521-samtrpti}



\begin{quote}
Before the new one: say the Malayalam for unease, then the Malayalam for to sob.
\end{quote}

\subsection*{You'll want to know: \ml{സംതൃപ്തി}}
\textbf{\ml{സംതൃപ്തി}} — \emph{saṁtṛpti} — \textquotedblleft{}satisfaction; contentment\textquotedblright{}.

\begin{quote}
\ml{ജോലി} \ml{നന്നായി} \ml{ചെയ്തതിൽ} \ml{അവന്} \ml{സംതൃപ്തി} \ml{തോന്നി}.
\end{quote}

\begin{quote}
\emph{jōli nannāyi ceytatil avanŭ saṁtṛpti tōnni.}
\end{quote}

\begin{quote}
He felt satisfied at having done the work well.
\end{quote}

Often with \ml{തോന്നുക} or \ml{ഉണ്ട്}: \ml{എനിക്ക്} \ml{സംതൃപ്തി} \ml{ഉണ്ട്}, I am content.

\begin{cousinweb}
\textbf{\ml{സംതൃപ്തി}} (\emph{saṁtṛpti}) is borrowed from Sanskrit: \emph{tṛpti}, \textquotedblleft{}satisfaction, contentment\textquotedblright{} (Monier-Williams), with the prefix \emph{saṁ-}, \textquotedblleft{}fully\textquotedblright{}.
\end{cousinweb}

\begin{grammarlens}[title={What you've built}]
One more word for describing what you lived and felt.
\end{grammarlens}

\subsection*{Your turn}
\begin{itemize}
\raggedright
  \item \emph{Say it:} \emph{saṁtṛpti}
  \item \emph{Say it:} \emph{saṁtṛpti}, once more
  \item \emph{Say it:} \emph{jōli nannāyi ceytatil avanŭ saṁtṛpti tōnni.} — He felt satisfied at having done the work well.
  \item \emph{Recall:} say the Malayalam for unease, then the Malayalam for to sob, then say \emph{saṁtṛpti} again
\end{itemize}

\begin{tcolorbox}[breakable,colback=teal!4,colframe=teal!35!black,title={Before you move on}]
What does \ml{സംതൃപ്തി} mean? (\textquotedblleft{}Satisfaction; contentment\textquotedblright{}.) Say it once more.
\end{tcolorbox}
\section[\texorpdfstring{\ml{വിഷാദം} (viṣādaṁ)}{വിഷാദം (viṣādaṁ)}]{\ml{വിഷാദം} (viṣādaṁ) — melancholy; depression}

\label{lesson:ML-C521-visadam}



\begin{quote}
Before the new one: say the Malayalam for to sob, then the Malayalam for satisfaction.
\end{quote}

\subsection*{You'll want to know: \ml{വിഷാദം}}
\textbf{\ml{വിഷാദം}} — \emph{viṣādaṁ} — \textquotedblleft{}melancholy; depression\textquotedblright{}.

\begin{quote}
\ml{മഴക്കാലത്ത്} \ml{ചിലർക്ക്} \ml{വിഷാദം} \ml{തോന്നാറുണ്ട്}.
\end{quote}

\begin{quote}
\emph{maḻakkālattŭ cilarkkŭ viṣādaṁ tōnnāṟuṇṭŭ.}
\end{quote}

\begin{quote}
Some people tend to feel low in the rainy season.
\end{quote}

In medicine \ml{വിഷാദരോഗം} is depression the illness. -\ml{ആറുണ്ട്} says something tends to happen.

\begin{cousinweb}
\textbf{\ml{വിഷാദം}} (\emph{viṣādaṁ}) is borrowed from Sanskrit \emph{viṣāda}, \textquotedblleft{}drooping, dejection, despondency\textquotedblright{} (Monier-Williams).
\end{cousinweb}

\begin{grammarlens}[title={What you've built}]
One more word for describing what you lived and felt.
\end{grammarlens}

\subsection*{Your turn}
\begin{itemize}
\raggedright
  \item \emph{Say it:} \emph{viṣādaṁ}
  \item \emph{Say it:} \emph{viṣādaṁ}, once more
  \item \emph{Say it:} \emph{maḻakkālattŭ cilarkkŭ viṣādaṁ tōnnāṟuṇṭŭ.} — Some people tend to feel low in the rainy season.
  \item \emph{Recall:} say the Malayalam for to sob, then the Malayalam for satisfaction, then say \emph{viṣādaṁ} again
\end{itemize}

\begin{tcolorbox}[breakable,colback=teal!4,colframe=teal!35!black,title={Before you move on}]
What does \ml{വിഷാദം} mean? (\textquotedblleft{}Melancholy; depression\textquotedblright{}.) Say it once more.
\end{tcolorbox}
\section[\texorpdfstring{\ml{പുഞ്ചിരിക്കുക} (puñcirikkuka)}{പുഞ്ചിരിക്കുക (puñcirikkuka)}]{\ml{പുഞ്ചിരിക്കുക} (puñcirikkuka) — to smile}

\label{lesson:ML-C521-puncirikkuka}



\begin{quote}
Before the new one: say the Malayalam for satisfaction, then the Malayalam for melancholy.
\end{quote}

\subsection*{You'll want to know: \ml{പുഞ്ചിരിക്കുക}}
\textbf{\ml{പുഞ്ചിരിക്കുക}} — \emph{puñcirikkuka} — \textquotedblleft{}to smile\textquotedblright{}.

\begin{quote}
\ml{ഫോട്ടോ} \ml{എടുത്തപ്പോൾ} \ml{എല്ലാവരും} \ml{പുഞ്ചിരിച്ചു}.
\end{quote}

\begin{quote}
\emph{phōṭṭō eṭuttappōḷ ellāvaruṁ puñciriccu.}
\end{quote}

\begin{quote}
Everyone smiled when the photo was taken.
\end{quote}

Past: \ml{പുഞ്ചിരിച്ചു}, \emph{puñciriccu}.

The noun is \ml{പുഞ്ചിരി}, a smile; \ml{ചിരിക്കുക} is to laugh.

\begin{cousinweb}
\textbf{\ml{പുഞ്ചിരിക്കുക}} (\emph{puñcirikkuka}) holds \ml{ചിരി}, \textquotedblleft{}laughter\textquotedblright{}, which is native Dravidian (DEDR 1562); where the first part comes from is not settled.
\end{cousinweb}

\begin{grammarlens}[title={What you've built}]
One more word for describing what you lived and felt.
\end{grammarlens}

\subsection*{Your turn}
\begin{itemize}
\raggedright
  \item \emph{Say it:} \emph{puñcirikkuka}
  \item \emph{Say it:} \emph{puñcirikkuka}, once more
  \item \emph{Say it:} \emph{phōṭṭō eṭuttappōḷ ellāvaruṁ puñciriccu.} — Everyone smiled when the photo was taken.
  \item \emph{Recall:} say the Malayalam for satisfaction, then the Malayalam for melancholy, then say \emph{puñcirikkuka} again
\end{itemize}

\begin{tcolorbox}[breakable,colback=teal!4,colframe=teal!35!black,title={Before you move on}]
What does \ml{പുഞ്ചിരിക്കുക} mean? (\textquotedblleft{}To smile\textquotedblright{}.) Say it once more.
\end{tcolorbox}
\section[\texorpdfstring{\ml{ഗൃഹാതുരത്വം} (gṛhāturatvaṁ)}{ഗൃഹാതുരത്വം (gṛhāturatvaṁ)}]{\ml{ഗൃഹാതുരത്വം} (gṛhāturatvaṁ) — nostalgia; homesickness}

\label{lesson:ML-C521-grhaturatvam}



\begin{quote}
Before the new one: say the Malayalam for melancholy, then the Malayalam for to smile.
\end{quote}

\subsection*{You'll want to know: \ml{ഗൃഹാതുരത്വം}}
\textbf{\ml{ഗൃഹാതുരത്വം}} — \emph{gṛhāturatvaṁ} — \textquotedblleft{}nostalgia; homesickness\textquotedblright{}.

\begin{quote}
\ml{പഴയ} \ml{പാട്ടുകൾ} \ml{കേൾക്കുമ്പോൾ} \ml{എനിക്ക്} \ml{ഗൃഹാതുരത്വം} \ml{തോന്നും}.
\end{quote}

\begin{quote}
\emph{paḻaya pāṭṭukaḷ kēḷkkumpōḷ enikkŭ gṛhāturatvaṁ tōnnuṁ.}
\end{quote}

\begin{quote}
When I hear old songs, I feel nostalgic.
\end{quote}

\ml{ഗൃഹം} house + \ml{ആതുരം} sick + -\ml{ത്വം} -ness: the state of being sick for home.

\begin{cousinweb}
\textbf{\ml{ഗൃഹാതുരത്വം}} (\emph{gṛhāturatvaṁ}) is a learned compound built from Sanskrit \emph{gṛha}, \textquotedblleft{}a house, a home\textquotedblright{}, and \emph{ātura}, \textquotedblleft{}sick in body or mind\textquotedblright{} (Monier-Williams), with the noun ending -\ml{ത്വം}.
\end{cousinweb}

\begin{grammarlens}[title={What you've built}]
One more word for describing what you lived and felt.
\end{grammarlens}

\subsection*{Your turn}
\begin{itemize}
\raggedright
  \item \emph{Say it:} \emph{gṛhāturatvaṁ}
  \item \emph{Say it:} \emph{gṛhāturatvaṁ}, once more
  \item \emph{Say it:} \emph{paḻaya pāṭṭukaḷ kēḷkkumpōḷ enikkŭ gṛhāturatvaṁ tōnnuṁ.} — When I hear old songs, I feel nostalgic.
  \item \emph{Recall:} say the Malayalam for melancholy, then the Malayalam for to smile, then say \emph{gṛhāturatvaṁ} again
\end{itemize}

\begin{tcolorbox}[breakable,colback=teal!4,colframe=teal!35!black,title={Before you move on}]
What does \ml{ഗൃഹാതുരത്വം} mean? (\textquotedblleft{}Nostalgia; homesickness\textquotedblright{}.) Say it once more.
\end{tcolorbox}
\section[\texorpdfstring{\ml{നെടുവീർപ്പിടുക} (neṭuvīrppiṭuka)}{നെടുവീർപ്പിടുക (neṭuvīrppiṭuka)}]{\ml{നെടുവീർപ്പിടുക} (neṭuvīrppiṭuka) — to sigh}

\label{lesson:ML-C521-netuvirppituka}



\begin{quote}
Before the new one: say the Malayalam for to smile, then the Malayalam for nostalgia.
\end{quote}

\subsection*{You'll want to know: \ml{നെടുവീർപ്പിടുക}}
\textbf{\ml{നെടുവീർപ്പിടുക}} — \emph{neṭuvīrppiṭuka} — \textquotedblleft{}to sigh\textquotedblright{}.

\begin{quote}
\ml{ബിൽ} \ml{കണ്ട്} \ml{അച്ഛൻ} \ml{നെടുവീർപ്പിട്ടു}.
\end{quote}

\begin{quote}
\emph{bil kaṇṭŭ acchan neṭuvīrppiṭṭu.}
\end{quote}

\begin{quote}
Father sighed when he saw the bill.
\end{quote}

Past: \ml{നെടുവീർപ്പിട്ടു}, \emph{neṭuvīrppiṭṭu}.

\ml{നെടുവീർപ്പ്}, a sigh, is a long breath; \ml{ഇടുക}, to put, makes it a verb.

\begin{cousinweb}
\textbf{\ml{നെടുവീർപ്പിടുക}} (\emph{neṭuvīrppiṭuka}) is a native Dravidian compound: \ml{നെടു}, \textquotedblleft{}long\textquotedblright{} (DEDR 3738), \ml{വീർപ്പ്}, \textquotedblleft{}breathing, a sigh\textquotedblright{} (DEDR 645), and \ml{ഇടുക}, \textquotedblleft{}to put\textquotedblright{} (DEDR 442).
\end{cousinweb}

\begin{grammarlens}[title={What you've built}]
One more word for describing what you lived and felt.
\end{grammarlens}

\subsection*{Your turn}
\begin{itemize}
\raggedright
  \item \emph{Say it:} \emph{neṭuvīrppiṭuka}
  \item \emph{Say it:} \emph{neṭuvīrppiṭuka}, once more
  \item \emph{Say it:} \emph{bil kaṇṭŭ acchan neṭuvīrppiṭṭu.} — Father sighed when he saw the bill.
  \item \emph{Recall:} say the Malayalam for to smile, then the Malayalam for nostalgia, then say \emph{neṭuvīrppiṭuka} again
\end{itemize}

\begin{tcolorbox}[breakable,colback=teal!4,colframe=teal!35!black,title={Before you move on}]
What does \ml{നെടുവീർപ്പിടുക} mean? (\textquotedblleft{}To sigh\textquotedblright{}.) Say it once more.
\end{tcolorbox}
